\begin{ex}
Số tròn trăm liền sau số 200 là số nào?
\shortans{300}
\loigiai{
Số tròn trăm liền sau 200 là 300.
}
\end{ex}

\begin{ex}
Số tròn trăm liền trước số 600 là số nào?
\shortans{500}
\loigiai{
Số tròn trăm liền trước 600 là 500.
}
\end{ex}

\begin{ex}
Số tròn trăm lớn nhất có ba chữ số là số nào?
\shortans{900}
\loigiai{
Số đó là 900.
}
\end{ex}

\begin{ex}
Số tròn trăm bé nhất có ba chữ số là số nào?
\shortans{100}
\loigiai{
Số đó là 100.
}
\end{ex}

\begin{ex}
Có bao nhiêu số tròn trăm có ba chữ số? (100, 200, ..., 900).
\shortans{9}
\loigiai{
Có đúng 9 số.
}
\end{ex}

\begin{ex}
Tính: $200 + 300$.
\shortans{500}
\loigiai{
Ta có $200 + 300 = 500$.
}
\end{ex}

\begin{ex}
Tính: $400 + 500$.
\shortans{900}
\loigiai{
Ta có $400 + 500 = 900$.
}
\end{ex}

\begin{ex}
Tính: $700 - 200$.
\shortans{500}
\loigiai{
Ta có $700 - 200 = 500$.
}
\end{ex}

\begin{ex}
Tính: $900 - 600$.
\shortans{300}
\loigiai{
Ta có $900 - 600 = 300$.
}
\end{ex}

\begin{ex}
Tính: $1000 - 400$.
\shortans{600}
\loigiai{
Ta có $1000 - 400 = 600$.
}
\end{ex}

\begin{ex}
Tính giá trị biểu thức: $300 + 400 - 200$.
\shortans{500}
\loigiai{
Ta có $300 + 400 = 700$, rồi $700 - 200 = 500$.
}
\end{ex}

\begin{ex}
Tính giá trị biểu thức: $800 - 500 + 400$.
\shortans{700}
\loigiai{
Ta có $800 - 500 = 300$, rồi $300 + 400 = 700$.
}
\end{ex}

\begin{ex}
Tính giá trị biểu thức: $1000 - 300 - 200$.
\shortans{500}
\loigiai{
Ta có $1000 - 300 = 700$, rồi $700 - 200 = 500$.
}
\end{ex}

\begin{ex}
Số tròn chục liền sau số 140 là số nào?
\shortans{150}
\loigiai{
Số tròn chục liền sau 140 là 150.
}
\end{ex}

\begin{ex}
Số tròn chục liền trước số 280 là số nào?
\shortans{270}
\loigiai{
Số tròn chục liền trước 280 là 270.
}
\end{ex}

\begin{ex}
Điền số thích hợp: $400 + \dots = 700$.
\shortans{300}
\loigiai{
Số cần điền: $700 - 400 = 300$.
}
\end{ex}

\begin{ex}
Điền số thích hợp: $800 - \dots = 300$.
\shortans{500}
\loigiai{
Số cần điền: $800 - 300 = 500$.
}
\end{ex}

\begin{ex}
Điền số thích hợp: $\dots + 500 = 1000$.
\shortans{500}
\loigiai{
Số cần điền: $1000 - 500 = 500$.
}
\end{ex}

\begin{ex}
Điền số thích hợp: $\dots - 200 = 600$.
\shortans{800}
\loigiai{
Số cần điền: $600 + 200 = 800$.
}
\end{ex}

\begin{ex}
Khối lớp Hai có 200 học sinh, khối lớp Ba có 300 học sinh. Cả hai khối có bao nhiêu học sinh?
\shortans{500}
\loigiai{
Cả hai khối có: $200 + 300 = 500$ học sinh.
}
\end{ex}

\begin{ex}
Một cửa hàng buổi sáng bán được 400 kg gạo, buổi chiều bán được 500 kg gạo. Cả ngày cửa hàng bán được bao nhiêu ki-lô-gam gạo?
\shortans{900}
\loigiai{
Cả ngày bán được: $400 + 500 = 900$ kg.
}
\end{ex}

\begin{ex}
Kho thóc có 800 bao thóc, người ta chuyển đi 300 bao thóc. Trong kho còn lại bao nhiêu bao thóc?
\shortans{500}
\loigiai{
Trong kho còn lại: $800 - 300 = 500$ bao.
}
\end{ex}

\begin{ex}
Vườn cây có 600 cây ăn quả, trong đó có 200 cây cam, còn lại là cây xoài. Vườn có bao nhiêu cây xoài?
\shortans{400}
\loigiai{
Số cây xoài: $600 - 200 = 400$ cây.
}
\end{ex}

\begin{ex}
Tổng của số tròn trăm lớn nhất có ba chữ số và số tròn trăm bé nhất có ba chữ số là bao nhiêu?
\shortans{1000}
\loigiai{
Tổng: $900 + 100 = 1000$.
}
\end{ex}

\begin{ex}
Hiệu của số tròn trăm lớn nhất có ba chữ số và số tròn trăm bé nhất có ba chữ số là bao nhiêu?
\shortans{800}
\loigiai{
Hiệu: $900 - 100 = 800$.
}
\end{ex}
